% Created 2026-04-27 lun 14:20
% Intended LaTeX compiler: pdflatex
\documentclass[11pt]{article}
\usepackage[utf8]{inputenc}
\usepackage[T1]{fontenc}
\usepackage{graphicx}
\usepackage{longtable}
\usepackage{wrapfig}
\usepackage{rotating}
\usepackage[normalem]{ulem}
\usepackage{amsmath}
\usepackage{amssymb}
\usepackage{capt-of}
\usepackage{hyperref}
\date{2026-04-27}
\title{Viaje con Inteligencia - Dossier Técnico y Estratégico}
\hypersetup{
 pdfauthor={M. Castillo},
 pdftitle={Viaje con Inteligencia - Dossier Técnico y Estratégico},
 pdfkeywords={},
 pdfsubject={},
 pdfcreator={Emacs 29.3 (Org mode 9.4.6)}, 
 pdflang={Spanish}}
\begin{document}

\maketitle
\tableofcontents


\section*{TABLE OF CONTENTS\hfill{}\textsc{no\_export}}
\label{sec:org53613ad}
\begin{itemize}
\item \hyperref[sec:org76da974]{Introducción}
\item \hyperref[sec:orgcb094a1]{Arquitectura del Sistema}
\item \hyperref[sec:org4c806f0]{Fuentes de Datos}
\item \hyperref[sec:orgb426034]{Inteligencia Artificial y Machine Learning}
\item \hyperref[sec:orgb572a20]{Diferenciación}
\item \hyperref[sec:org3e74d2b]{Canal de Telegram}
\item \hyperref[sec:org7f38bb1]{Integración Mastodon}
\item \hyperref[sec:org076d7f8]{Modelo de Negocio}
\item \hyperref[sec:org0d514bc]{Potencial de Crecimiento}
\item \hyperref[sec:org28fae3a]{Tecnología}
\item \hyperref[sec:org246f1d0]{Roadmap}
\item \hyperref[sec:org902dbd2]{Anexos}
\end{itemize}

\section{INTRODUCCIÓN\hfill{}\textsc{no\_export}}
\label{sec:org76da974}
\subsection{Visión del Proyecto}
\label{sec:orgdd6be1e}

Viaje con Inteligencia es una plataforma de planificación de viajes que utiliza inteligencia artificial y machine learning para ofrecer recomendaciones personalizadas, alertas en tiempo real y análisis que ningún otro servicio proporciona. A diferencia de las aplicaciones de viaje tradicionales que operan con datos estáticos y recomendaciones genéricas, nuestro sistema aprende de las preferencias del usuario, integra fuentes oficiales en tiempo real y proporciona recomendaciones contextuales únicas.

\subsection{Problema que Resuelve}
\label{sec:orga0928bf}

El viajero moderno se enfrenta a un ecosistema fragmentado:
\begin{itemize}
\item Aplicaciones de地名 con información desactualizada
\item Guías de viaje genéricas que ignoran preferencias individuales
\item Alertas de seguridad dispersas en múltiples fuentes
\item cero personalización contextual
\end{itemize}

Viaje con Inteligencia unifica estas funciones en un solo ecosistema inteligente.

\subsection{Propuesta de Valor}
\label{sec:orgddc39ea}

\begin{center}
\begin{tabular}{lll}
Aspecto & Tradicional & Nuestra Solución\\
\hline
Datos & Estáticos, desactualizados & Tiempo real, oficiales\\
Recomendaciones & Genéricas & ML personalizado\\
Alertas & Fragmentadas & Unificadas, proactivas\\
Experiencia & Transaccional & Contextual, aprend\\
\end{tabular}
\end{center}

\subsection{Métricas Actuales (2026-04-27)}
\label{sec:org5c71dd6}

\begin{center}
\begin{tabular}{lr}
Métrica & Valor\\
\hline
Países cubiertos & 95+\\
Posts blog & 52\\
Usuarios Telegram & \textasciitilde{}50\\
Suscriptores newsletter & 12+\\
Rutas temáticas ML & 4\\
Regiones vinícolas & 8\\
\end{tabular}
\end{center}

\section{ARQUITECTURA DEL SISTEMA\hfill{}\textsc{no\_export}}
\label{sec:orgcb094a1}
\subsection{DIAGRAMAS ARQUITECTURA}
\label{sec:org619ca2d}

\lstset{language=text,label= ,caption= ,captionpos=b,numbers=none}
\begin{lstlisting}
[Usuario] --> [Frontend Next.js]
[Frontend] --> [API Routes]
[API Routes] --> [Supabase]
[API Routes] --> [Telegram Bot]
[API Routes] --> [Mastodon]
[API Routes] --> [Resend]
[ML Engine] --> [Clustering]
[ML Engine] --> [Recommendations]
[INE/MAEC/OMS] --> [API Routes]
\end{lstlisting}

\subsection{Componentes Principales}
\label{sec:org95b90ff}

\begin{center}
\begin{tabular}{lll}
Componente & Tecnología & Función\\
\hline
Frontend & Next.js 14 & UI/reactiva\\
Backend & Next.js API Routes & Endpoints serverless\\
Base de datos & Supabase & PostgreSQL + Auth\\
ML Engine & Python/TypeScript & Algoritmos\\
Telegram Bot & Telegram API & Interacciones\\
Email & Resend & Newsletters\\
Hosting & Vercel & Despliegue\\
\end{tabular}
\end{center}

\subsection{APIs Principales}
\label{sec:orgf3af991}

\begin{verbatim}
[Frontend] ---> [API Routes] ---> [Supabase]
                    |
                    +--> [Telegram Bot]
                    |
                    +--> [Mastodon]
                    |
                    +--> [Resend]

[INE/MAEC/OMS] ---> [API Routes]
\end{verbatim}

\section{FUENTES DE DATOS\hfill{}\textsc{no\_export}}
\label{sec:org4c806f0}
\subsection{Fuentes Oficiales Integradas}
\label{sec:org34948df}

\begin{center}
\begin{tabular}{llll}
Fuente & Datos & Frecuencia & Uso\\
\hline
MAEC & Alertas viaje & Diario & Scraper riesgo\\
INE & Turistas, occupancy & Mensual & IST, ML\\
OMS & Alertas sanitarias & Variable & Salud viajeros\\
AEMET & Meteorología & Hora & Clima tiempo real\\
UNWTO & Estadísticas globales & Anual & KPIs turismo\\
\end{tabular}
\end{center}

\subsection{Detalle de Fuentes}
\label{sec:org65a2589}

\subsubsection{MAEC - Ministerio de Asuntos Exteriores}
\label{sec:org7c8d67d}

El sistema scraping del MAEC se ejecuta diariamente a las 06:00 UTC, extrayendo:
\begin{itemize}
\item Nivel de riesgo por país
\item Recomendaciones específicas
\item Enlaces a fichas país PDF
\item Últimas alertas de seguridad
\end{itemize}

\subsubsection{Flujo Scraping MAEC}
\label{sec:org9973216}

\begin{verbatim}
[Init] --> [06:00 UTC Daily]
[06:00 UTC Daily] --> [Fetch MAEC alerts]
[Fetch MAEC alerts] --> [Parse HTML]
[Parse HTML] --> [Store Supabase]
[Store Supabase] --> [Notify Telegram/Email]
[Notify Telegram/Email] --> [Done]
\end{verbatim}

\subsubsection{INE - Instituto Nacional de Estadística}
\label{sec:orgc96cf3b}

Datos de turismo de España (FRONTUR/EGATUR):
\begin{itemize}
\item Llegadas de turistas internacionales
\item Pernoctaciones por origen
\item Estancia media
\item Ingresos por turismo
\end{itemize}

\subsection{Sistema de Alertas}
\label{sec:org33cb02c}

\begin{center}
\begin{tabular}{lll}
Tipo & Canal & Frecuencia\\
\hline
Riesgo país & Telegram Bot & On-demand\\
Resumen diario & Email + Telegram & 00:00 UTC\\
Digest semanal & Email + Telegram + Mastodon & Lunes 08:00\\
\end{tabular}
\end{center}

\section{INTELIGENCIA ARTIFICIAL Y MACHINE LEARNING\hfill{}\textsc{no\_export}}
\label{sec:orgb426034}
\subsection{Stack de ML Implementado}
\label{sec:org5473f64}

\begin{center}
\begin{tabular}{lll}
Algoritmo & Uso & Datos\\
\hline
K-Means Clustering & Segmentación destinos & Riesgo, IPC, tourist arrivals\\
Preference Matching & Recomendaciones & user prefs + destination attributes\\
Risk Scoring & Análisis riesgo & MAEC + OMS + historical\\
Time Series & Predicción precios & Historical INE\\
\end{tabular}
\end{center}

\subsection{Sistema de Recomendaciones}
\label{sec:orgd82b968}

\begin{verbatim}
Input:  [User Prefs] --> [Budget] --> [Dates]
                         |           |
                         v           v
Feature Engineering --> Model Scoring --> Ranking --> Output (Top 3)

Scoring: (prefMatch × 3) + (value × 2) + (safety × 1.5)
\end{verbatim}

\subsubsection{Algoritmo de Scoring}
\label{sec:org7e757d1}

@enduml
\#+end\textsubscript{src}

\subsection{Scoring Algorithm}
\label{sec:orgbe5041b}

\begin{verbatim}
preferenceMatchScore = Σ(pref_weight × attribute_match)
valueScore = (price_historical / current_price) × 100
safetyScore = (10 - risk_level) × 10
crowdScore = ist_index × season_multiplier

finalScore = (preferenceMatchScore × 0.4) + 
           (valueScore × 0.25) + 
           (safetyScore × 0.20) + 
           (crowdScore × 0.15)
\end{verbatim}

\subsection{ML Features por Ruta}
\label{sec:orgd35ecb0}

\begin{center}
\begin{tabular}{lrrrl}
Ruta & Popularity & Safety & Value & Best Season\\
\hline
Molinos La Mancha & 7.5 & 9.5 & 8.5 & Abr-May, Sept-Oct\\
Faros Costa & 8.5 & 9.0 & 7.0 & May-Jun, Sept-Oct\\
Murcia Interior & 6.0 & 9.5 & 9.0 & Mar-May, Oct-Nov\\
Rutas del Vino & 7.0 & 9.5 & 8.5 & Mar-May, Sept-Nov\\
\end{tabular}
\end{center}

\subsection{clustering de Destinos}
\label{sec:orgb5e309e}

El sistema utiliza K-Means clustering con 7 segmentos:
\begin{enumerate}
\item 🏖️ Playa - destinos costeros
\item 🏛️ Cultural - patrimonio histórico
\item 👨‍👩‍👧‍👦 Familiar - seguro, económico
\item 🏔️ Naturaleza - aventura, naturaleza
\item 💼 Negocio - centros business
\item Salud - wellbeing, spa
\item 🏔️ Montana - ski, montaña
\end{enumerate}

\section{DIFERENCIACIÓN\hfill{}\textsc{no\_export}}
\label{sec:orgb572a20}
\subsection{Comparativa con Competidores}
\label{sec:orgceb8615}

\begin{center}
\begin{tabular}{lllll}
Aspecto & TripAdvisor & Lonely Planet & Booking & Viaje con IA\\
\hline
Datos oficiales & ❌ Parcial & ✅ Estático & ❌ & ✅ Tiempo real\\
ML personalized & ❌ & ❌ & ❌ & ✅\\
Alertas MAEC & ❌ & ❌ & ❌ & ✅\\
Rutas temáticas ML & ❌ & ❌ & ❌ & ✅\\
Telegram Bot & ❌ & ❌ & ❌ & ✅\\
Mastodon & ❌ & ❌ & ❌ & ✅\\
\end{tabular}
\end{center}

\subsection{Ventajas Competitivas}
\label{sec:orgd0e07fa}

\subsubsection{1. Datos Oficiales, Tiempo Real}
\label{sec:org416fa64}

比其他平台更新的数据，来自官方来源：
\begin{itemize}
\item MAEC alerts → 每日更新
\item INE statistics → 每月更新
\item OMS health → 实时更新
\end{itemize}

\subsubsection{2. Machine Learning Personalizado}
\label{sec:orgc3e0df2}

我们的推荐引擎不是"如果-那么"规则，而是从以下学习：
\begin{itemize}
\item 类似用户的行为模式
\item 您明确的偏好
\item 隐式反馈（点击、停留时间）
\end{itemize}

\subsubsection{3. Integración Completa}
\label{sec:org50b9888}

单一平台提供：
\begin{itemize}
\item 搜索 → 推荐 → 预订全部
\item 警报 → 风险分析
\item 行程 → 优化
\end{itemize}

\subsubsection{4. Canales Múltiples}
\label{sec:org96a7a2a}

\begin{center}
\begin{tabular}{ll}
Canal & Función\\
\hline
Web & Búsqueda, recomendaciones\\
Telegram Bot & Alertas interactivas\\
Telegram Channel & Newsletter automático\\
Mastodon & Presencia social\\
Email & Resumen semanal\\
\end{tabular}
\end{center}

\subsection{Propuesta Única de Venta (PUV)}
\label{sec:org44f5b8a}

> "La única plataforma de viajes en español que combina inteligencia artificial + machine learning con datos oficiales en tiempo real para ofrecer recomendaciones verdaderamente personalizadas."

\section{CANAL DE TELEGRAM\hfill{}\textsc{no\_export}}
\label{sec:org3e74d2b}
\subsection{Arquitectura Telegram}
\label{sec:org96f6e3e}

\begin{verbatim}
[Canal: @ViajeConInteligencia] <-- [Bot interactivo]
[Bot: @ViajeConInteligenciaBot] <-- [API Routes]

Backend --> [sendTelegramMessage] --> [Canal]
         --> [cron/digest] --> [Canal]
         --> [newsletter] --> [Canal]

Funciones: Alertas riesgo + Resumen semanal + Posts blog
\end{verbatim}

\subsection{Comandos Disponibles}
\label{sec:orge6b172f}
[Alertas riesgo] ..> [Bot: @ViajeConInteligenciaBot]
[Resumen semanal] ..> [Canal: @ViajeConInteligencia]
[Posts blog] ..> [Canal: @ViajeConInteligencia]

note right of [Canal: @ViajeConInteligencia]
  Recibe posts automáticamente
  cada semana + alertas
end note

@enduml
\#+end\textsubscript{src}

\subsection{Comandos Disponibles}
\label{sec:org2e1522c}

\begin{center}
\begin{tabular}{ll}
Comando & Función\\
\hline
/start & Bienvenida + información\\
/alertasviaje & Resumen de alertas de riesgo\\
/alertasviaje full & Detalle completo\\
/help & Ayuda\\
/suscribirse & Alta en newsletter\\
\end{tabular}
\end{center}

\subsection{Integración con ML}
\label{sec:org63df571}

El Bot utiliza el sistema de recomendaciones ML para:
\begin{itemize}
\item Responder consultas de países específicos
\item Proporcionar recomendaciones personalizadas
\item Analizar riesgo contextual
\end{itemize}

\subsection{Métricas de Uso}
\label{sec:orgebd2cd5}

\begin{center}
\begin{tabular}{ll}
Métrica & Valor\\
\hline
Usuarios activos & \textasciitilde{}50\\
Mensajes/día & Variable\\
Suscriptores canal & Incrementando\\
\end{tabular}
\end{center}

\section{INTEGRACIÓN MASTODON\hfill{}\textsc{no\_export}}
\label{sec:org7f38bb1}
\subsection{Presencia en Mastodon}
\label{sec:orgc340f90}

\begin{center}
\begin{tabular}{ll}
Aspecto & Valor\\
\hline
Cuenta & @viajeinteligencia@mastodon.social\\
Seguidores & Growing\\
Publicaciones & Automáticas (cron) + Manuales\\
\end{tabular}
\end{center}

\subsection{Sistema de Publicación Automática}
\label{sec:orgf765a69}

\begin{center}
\begin{tabular}{ll}
Trigger & Contenido\\
\hline
Nuevo post blog & Resumen + hashtags\\
Digest semanal & Top posts + alertas\\
Alerta MAEC importante & Enlace a análisis\\
\end{tabular}
\end{center}

\subsection{API de Publicación}
\label{sec:org50d761a}

\begin{verbatim}
Triggers:Nuevo blog post | cron/digest | manual
API: POST /api/v1/statuses
Token: MASTODON_ACCESS_TOKEN
Visibility: public
\end{verbatim}

\section{MODELO DE NEGOCIO\hfill{}\textsc{no\_export}}
\label{sec:org076d7f8}
\subsection{Estructura de Ingresos}
\label{sec:org50cb266}

\begin{center}
\begin{tabular}{lll}
Tier & Precio & Funcionalidades\\
\hline
Free & 0€ & Búsqueda país, alertas básicas, 3 posts blog/mes\\
Premium & 4.99€/mes & ML recomendaciones completas, newsletter, PDF checklist\\
yearly & 19.99€/año & Todo premium + priority alerts\\
\end{tabular}
\end{center}

\subsection{Funcionalidades por Tier}
\label{sec:org260d364}

\begin{center}
\begin{tabular}{lll}
Feature & Free & Premium\\
\hline
Países en mapa & ✅ & ✅\\
Alertas MAEC & Básica & Completa\\
ML recomendaciones & 3/día & Ilimitado\\
Newsletter & - & ✅\\
PDF checklist & - & ✅\\
Priority alerts & - & ✅\\
Historial favoritos & - & ✅\\
\end{tabular}
\end{center}

\subsection{Suscripciones (Stripe)}
\label{sec:orgd21b605}

\begin{center}
\begin{tabular}{lll}
Métrica & Actual & Target Q2 2026\\
\hline
Suscripciones & 0 & 100\\
MRR & 0€ & \textasciitilde{}300€\\
Churn & - & <5\%\\
\end{tabular}
\end{center}

\section{POTENCIAL DE CRECIMIENTO\hfill{}\textsc{no\_export}}
\label{sec:org0d514bc}
\subsection{Análisis de Mercado}
\label{sec:org2378758}

\begin{center}
\begin{tabular}{lll}
Segmento & TAM & Share Viabilidad\\
\hline
Viajes España & 100M€/año & 0.1\% → 1\% = 1M€\\
Viajes outbound España & 30M personas/año & 0.1\% → 0.5\% = 150K\\
Enoturismo & 5M€/año & 1\% → 5\% = 250K€\\
\end{tabular}
\end{center}

\subsection{Factores de Crecimiento}
\label{sec:org97bca0a}

\subsubsection{1. SEO y Contenido}
\label{sec:orgd6c35bd}
\begin{itemize}
\item 52+ posts blog indexados
\item Keywords long-tail: "es seguro viajar a [país]"
\item Datos estructurados JSON-LD
\end{itemize}

\subsubsection{2. Red de Distribución}
\label{sec:orgeaeb569}
\begin{itemize}
\item Telegram: \textasciitilde{}50 usuarios base
\item Mastodon: presencia social
\item Email: 12+ suscriptores
\end{itemize}

\subsubsection{3. Viralidad}
\label{sec:org8350b1a}
\begin{itemize}
\item Recomendaciones ML comparten
\item Alertas de riesgo comparten
\end{itemize}

\subsection{Proyección 12 Meses}
\label{sec:org06743dd}

\begin{center}
\begin{tabular}{llll}
Métrica & Q2 2026 & Q3 2026 & Q4 2026\\
\hline
Visitantes/mes & 1,000 & 5,000 & 10,000\\
Usuarios registrados & 100 & 500 & 1,000\\
Suscripciones premium & 25 & 100 & 250\\
Ingresos MRR & 100€ & 400€ & 1,000€\\
\end{tabular}
\end{center}

\subsection{Expansión Geográfica}
\label{sec:org9aaa7a6}

\begin{center}
\begin{tabular}{rll}
Fase & Mercados & Timeline\\
\hline
1 & España (base) & Done\\
2 & Latinoamérica hispanohablante & Q3 2026\\
3 & Portugal & Q4 2026\\
4 & English market & 2027\\
\end{tabular}
\end{center}

\section{TECNOLOGÍA\hfill{}\textsc{no\_export}}
\label{sec:org28fae3a}
\subsection{Stack Técnico}
\label{sec:org15b99d1}

\begin{center}
\begin{tabular}{ll}
Capa & Tecnología\\
\hline
Frontend & Next.js 14, React, Tailwind\\
Backend & Next.js API Routes\\
Base de datos & Supabase (PostgreSQL)\\
Auth & Supabase Auth\\
ML & TypeScript + Python\\
Email & Resend\\
Pagos & Stripe\\
Hosting & Vercel\\
CDN & Vercel Edge\\
\end{tabular}
\end{center}

\subsection{Dependencias Principales}
\label{sec:org242b033}

\begin{verbatim}
dependencies:
  next: ^14.0.0
  react: ^18.0.0
  @supabase/supabase-js: ^2.0.0
  @stripe/stripe-js: ^2.0.0
  react-leaflet: ^4.0.0
  dexie: ^4.0.0 (IndexedDB)
  jspdf: ^2.0.0
  resend: ^3.0.0
\end{verbatim}

\subsection{Infraestructura Vercel}
\label{sec:orga0153a5}

\begin{center}
\begin{tabular}{ll}
Recurso & Configuración\\
\hline
Functions & Serverless\\
Cron Jobs & 4 configured\\
Edge Functions & En desarrollo\\
Preview & PR automáticas\\
\end{tabular}
\end{center}

\subsection{Seguridad}
\label{sec:org4e291c5}

\begin{itemize}
\item Variables de entorno en Vercel
\item Auth con Supabase Auth
\item Rate limiting básico
\item CSRF protection
\item Sanitización de inputs
\end{itemize}

\section{ROADMAP\hfill{}\textsc{no\_export}}
\label{sec:org246f1d0}
\subsection{2026 Roadmap}
\label{sec:orgb4ec0a3}

\begin{center}
\begin{tabular}{llll}
Prioridad & Feature & Estado & Q\\
\hline
🔴 Alta & Programa afiliados & Pending & Q2\\
🟠 Media & AviationStack API & Pending & Q2\\
🟠 Media & Predicción precios ML & Pending & Q3\\
🟡 Baja & App móvil PWA & In progress & Q2\\
🟡 Baja & Anomaly detection & Pending & Q4\\
\end{tabular}
\end{center}

\subsection{Features Futuros}
\label{sec:org6fd5b87}

\begin{center}
\begin{tabular}{lll}
Feature & Descripción & Impacto\\
\hline
Predicción precios & ML predicciones vuelos & Alto\\
Recomendaciones predictivas & Sugerir antes de buscar & Alto\\
Integración calendario & Alarmas contextuales & Medio\\
Marketplace & Affiliate hoteless & Alto\\
\end{tabular}
\end{center}

\section{ANEXOS\hfill{}\textsc{no\_export}}
\label{sec:org902dbd2}
\subsection{URLs del Proyecto}
\label{sec:org22a458c}

\begin{center}
\begin{tabular}{ll}
Recurso & URL\\
\hline
Web & \url{https://viajeinteligencia.com}\\
Blog & \url{https://viajeinteligencia.com/blog}\\
Vercel & \url{https://viaje-con-inteligencia.vercel.app}\\
KPIs Paz & \url{https://viaje-con-inteligencia.vercel.app/kpi}\\
KPIs Turismo & \url{https://viaje-con-inteligencia.vercel.app/turismo}\\
Bot Telegram & \url{https://t.me/ViajeConInteligenciaBot}\\
Canal Telegram & @ViajeConInteligencia\\
Mastodon & @viajeinteligencia@mastodon.social\\
Repo GitHub & \url{https://github.com/mcasrom/viaje-con-inteligencia}\\
\end{tabular}
\end{center}

\subsection{Variables de Entorno}
\label{sec:org545d2c7}

\begin{center}
\begin{tabular}{ll}
Variable & Descripción\\
\hline
SUPABASE\textsubscript{URL} & URL proyecto Supabase\\
SUPABASE\textsubscript{ANON}\textsubscript{KEY} & Key pública Supabase\\
TELEGRAM\textsubscript{BOT}\textsubscript{TOKEN} & Token Bot Telegram\\
TELEGRAM\textsubscript{CHANNEL}\textsubscript{ID} & ID Canal Telegram\\
MASTODON\textsubscript{ACCESS}\textsubscript{TOKEN} & Token Mastodon\\
CRON\textsubscript{SECRET} & Secret para cron jobs\\
RESEND\textsubscript{API}\textsubscript{KEY} & Key API Resend\\
GROQ\textsubscript{API}\textsubscript{KEY} & Key API Groq (IA)\\
STRIPE\textsubscript{SECRET}\textsubscript{KEY} & Key secret Stripe\\
\end{tabular}
\end{center}

\subsection{Cron Jobs Configurados}
\label{sec:orgffcf73c}

\begin{center}
\begin{tabular}{lll}
Job & Schedule & Función\\
\hline
scrape-maec & 0 6 * * * & Scraping MAEC\\
check-alerts & 0 8 * * * & Verificar alertas\\
weekly-digest & 0 8 * * 1 & Digest semanal\\
daily-digest & 0 0 * * * & Resumen diario\\
\end{tabular}
\end{center}

\subsection{Changelog Reciente}
\label{sec:orgb63290f}

\begin{center}
\begin{tabular}{rl}
Fecha & Cambio\\
\hline
2026-04-27 & Blog post IA/ML Telegram + posts automáticos a Canal\\
2026-04-27 & Rutas del Vino (8 regiones) + ML optimization\\
2026-04-26 & INE Tourism ML clustering\\
2026-04-25 & Planificador Simple + preferencias ML\\
2026-04-24 & Reclamaciones PDF\\
2026-04-23 & PDF Premium completo\\
\end{tabular}
\end{center}

\subsection{Licencia y Contribución}
\label{sec:org80be6ab}

MIT License - contribuciones bienvenidas via GitHub.

\subsection{Contacto}
\label{sec:org82e12bc}

\begin{itemize}
\item Email: mybloggingnotes@gmail.com
\item Autor: M. Castillo
\item Fecha creación: 2024
\item Última actualización: 2026-04-27
\end{itemize}

\section{CONCLUSIÓN\hfill{}\textsc{no\_export}}
\label{sec:org277eef6}

Viaje con Inteligencia representa un cambio de paradigma en la planificación de viajes: de información estática y genérica a inteligencia contextual y personalizada. Con datos oficiales en tiempo real, machine learning avanzada y múltiples canales de distribución, la plataforma está posicionada para capturar el creciente mercado de viajeros españoles que buscan recomendaciones inteligentes.

El potencial de crecimiento es significativo, con múltiples vectores de expansión:
\begin{enumerate}
\item Más usuarios → más datos → mejores recomendaciones (efecto red)
\item Suscripciones premium → ingresos recurrentes
\item Programa afiliados → monetización adicional
\item Expansión geográfica → nuevos mercados
\end{enumerate}

La diferenciación clara vs. competidores tradicionales y el enfoque en IA/ML proporcionan una barrera competitiva sostenible.

---

\textbf{Documento generado: 2026-04-27}
\textbf{Versión: 1.0}
\textbf{Palabras: \textasciitilde{}2,800}
\end{document}